\section{Causal pair coupling at matched query counts}
\label{supp:cv-joint-selector}
To test whether the wider menu can be used causally, we fit fresh donor
forecasters inside the nine original fitting sources. Each of the
\CVSelectorDonor{} donors excludes its target source groups and every A
source. The shorter donor budget (\CVSelectorDonorSteps{} total updates)
does not match the maturity of the frozen A models. The latter are not
refit. Two cost heads per seed receive \CVSelectorHeadSteps{} total updates.
For pair edges, the same predictor scores the joint edit and its two
one-sided edits. The additive and coupled decoders share the one-sided
scores and paths; only the latter uses the difference between joint
and additive scores as a selection term. No A labels enter these choices.

At \CVSelectorEdits{} matched edits per arm over three seeds, the direct
coupled-minus-additive agent full-grid lower gain is \CVSelectorAgent{}
CV-normalized points (95\% source-bootstrap interval). The corresponding
paired-displacement lower gain is \CVSelectorPair{} pixels, despite an
observed-only gain of \CVSelectorObservedPair{} pixels. Matching
each query's total peak displacement without changing recipients gives
\CVSelectorMatchedAgent{} points and \CVSelectorMatchedPair{} pixels.
The agent lower contrast against independent ranking is
\CVSelectorIndependentAgent{} points and against a common candidate
family is \CVSelectorUniformAgent{} points.

Original CV-relative risk passes are \CVSelectorJointRisk{}/102 for
coupled, \CVSelectorAddRisk{}/102 for additive, and
\CVSelectorIndependentRisk{}/102 for independent. The exact-count
arms are forced comparisons, not operating policies with learned
admission. Pair-displacement error is not a physical safety measure.
These 34 exposed recordings and 14 source groups are not independent
confirmation, and the original threshold and support rules are unchanged.
